\chapter{چارچوب پیشنهاد فنی}
\begin{DSCBox}{راهنمای استفاده از این بخش}
محتوای این فصل نمونه است. عنوان‌ها و فیلدهای داخل کروشه را متناسب با پروژه تکمیل کنید؛ ردیف‌های جدول، تعهد اجرایی یا زمان‌بندی قطعی محسوب نمی‌شوند.
\end{DSCBox}
\section{خلاصه مدیریتی}
[مسئله سازمان، نتیجه مورد انتظار و ارزش پیشنهادی راهکار را در یک یا دو بند کوتاه بنویسید. مخاطب این بخش باید بتواند بدون ورود به جزئیات فنی، هدف همکاری را درک کند.]
\section{دامنه و خروجی‌ها}
\begin{itemize}
\item [خروجی اول و معیار پذیرش آن]
\item [خروجی دوم و معیار پذیرش آن]
\item [محدودیت‌ها، پیش‌نیازها و موارد خارج از دامنه]
\end{itemize}
\section{مراحل اجرا}
\renewcommand{\arraystretch}{1.35}
\begin{tabularx}{\linewidth}{@{}>{\raggedleft\arraybackslash}p{27mm}>{\raggedleft\arraybackslash}X>{\raggedleft\arraybackslash}p{43mm}@{}}
\rowcolor{DSCNavy}\textcolor{white}{\bfseries مرحله} & \textcolor{white}{\bfseries فعالیت} & \textcolor{white}{\bfseries خروجی}\\
\rowcolor{DSCPale}شناخت & تحلیل نیازها و شرایط موجود & سند نیازمندی‌ها\\
طراحی & بررسی معماری و گزینه‌های فنی & طرح فنی و برنامه ارزیابی\\
\rowcolor{DSCPale}نمونه‌سازی & پیاده‌سازی نمونه مورد توافق & نمونه قابل آزمون\\
ارزیابی & اجرای آزمون‌ها و تحلیل نتایج & گزارش ارزیابی و تحویل\\
\end{tabularx}
\section{معیارهای پذیرش}
[معیار، روش اندازه‌گیری، مقدار هدف و مسئول تأیید هر خروجی را مشخص کنید.]
\subsection{فرض‌ها و وابستگی‌ها}
[دسترسی‌ها، منابع پردازشی، داده‌ها و همکاری‌های موردنیاز را ثبت کنید.]
